\section{Public Cohort Results}
\label{sec:public-results}

The frozen public cohort completed all 144 planned rows with no invalid
infrastructure run. Table~\ref{tab:public-success} reports task success across
36 executions per condition. The preregistered primary comparison aggregates
the three repetitions within each of the 12 cases: canonical \memorix{} reached
97.2\% versus 94.4\% for no memory, an absolute difference of $+2.8$ points
(95\% case-clustered bootstrap interval $[-8.3, 16.7]$; descriptive sign-flip
$p=1.0$). The interval includes both a modest disadvantage and a meaningful
advantage, so this cohort does not support a stable success-rate claim.

\begin{table}[t]
\centering
\small
\caption{Public fixed-model cohort outcomes. Baseline rows are descriptive;
only Memorix versus no memory is the frozen primary comparison.}
\label{tab:public-success}
\begin{tabular}{lrr}
\toprule
Condition & Successful runs & Rate \\
\midrule
No memory & 34 / 36 & 94.4\% \\
Mem0 & 34 / 36 & 94.4\% \\
AgentMemory & 35 / 36 & 97.2\% \\
Canonical \memorix{} & 35 / 36 & 97.2\% \\
\bottomrule
\end{tabular}
\end{table}

Canonical \memorix{} used 19,479 input tokens per case-cluster on average,
compared with 23,360 for no memory (difference $-3{,}881$; 95\% interval
$[-20{,}008, 6{,}938]$). It also averaged 26.65 versus 28.62 seconds,
\$0.00402 versus \$0.00589 provider cost, and 8.17 versus 8.92 tool calls.
Every resource interval spans zero. These directionally favorable averages are
therefore efficiency observations, not evidence of a reliable reduction.

All six failed tasks exhausted the fixed 24-tool-call limit rather than failing
a verifier. The public failure patterns are informative: in
\texttt{python-header-policy-v1}, canonical \memorix{} succeeded in all three
repetitions while no memory succeeded once; in \texttt{go-slug-v1}, canonical
\memorix{} exhausted the limit once while no memory succeeded in all three.
The opposing examples explain the near-zero aggregate result and motivate a
larger real-repository cohort with explicitly varied dependency strength.

We then froze a separate, post-result replication under
\texttt{deepseek/deepseek-v4-flash}. It retained the same 12 cases, three
repetitions, exact tool policy, and canonical \memorix{} versus no-memory
contrast, but did not rerun secondary baseline systems. All 72 planned rows
were valid and all passed their behavioral verifier: both conditions therefore
had 36 / 36 successes and a zero success-rate difference. This is not evidence
that memory failed generally; rather, this strong model had no accuracy headroom
on these fully public fixtures.

The replication nevertheless exposes a cost boundary. Canonical \memorix{}
averaged 11,398 versus 9,192 input tokens and \$0.00111 versus \$0.00094
provider cost per case-cluster, differences of $+2{,}206$ tokens and
\$0.000174 respectively. Their descriptive 95\% cluster intervals excluded
zero. Wall time increased by 3.29 seconds on average, but its interval
$[-1.48, 8.10]$ spans zero. The model-specific contrast reinforces the intended
claim boundary: memory should earn its context and latency on tasks with a
real dependency, not be injected automatically because a prior session exists.
